\subsection{The group of classes of invertible modules}

*(We preserve the notation of nos. 5 and 6.)*

Under multiplication, the sub-A-modules of B form a commutative monoid M, with A as identity element. Then the invertible modules are the invertible elements of M and therefore form a commutative group J. We have seen (no. 6, Proposition 10) that the inverse of \(M \in J\) is A: M.

Let A* (resp. B*) be the multiplicative group of invertible elements of A (resp. B) and let u denote the canonical injection A → B. For all \(b \in B^{*}\) , \(\theta(b) = bA\) is an invertible sub-A-module. The mapping \(\theta: B^{*} \to J\) is a homomorphism whose kernel is \(u(A^{*})\) ; its cokernel will be denoted by C or C(A). The group C is called the group of classes of invertible sub-A-modules of B. The following exact sequence has been constructed

\[
(1) \longrightarrow A ^ {*} \stackrel {u} {\longrightarrow} B ^ {*} \stackrel {\theta} {\longrightarrow} \mathfrak {J} \stackrel {\rho} {\longrightarrow} \mathfrak {C} \longrightarrow (1)\tag{10}
\]

where (1) denotes the group consisting only of the identity element and \(\rho\) is the canonical mapping \(\mathfrak{J} \to \mathfrak{C} = \mathfrak{J}/\theta(\mathrm{B}^*)\) .

As every invertible sub-A-module M of B is projective of rank 1 (no. 6, Theorem 4), the element \(\operatorname{cl}(M) \in P(A)\) is defined (no. 4).

PROPOSITION 12. The mapping cl: \(\mathfrak{J} \to P(A)\) defines, by taking quotients, an isomorphism from \(\mathfrak{C} = \mathfrak{J}/\theta(B^{*})\) onto the kernel of the canonical homomorphism

\[
\phi \colon P (\mathrm{A}) \to P (\mathrm{B})
\]

(no. 4).

In other words, there is an exact sequence

\[
(1) \longrightarrow \mathrm{A} ^ {*} \stackrel {u} {\longrightarrow} \mathrm{B} ^ {*} \stackrel {\theta} {\longrightarrow} \mathfrak {J} \stackrel {\mathrm{cl}} {\longrightarrow} P (\mathrm{A}) \stackrel {\phi} {\longrightarrow} P (\mathrm{B}).\tag{11}
\]

It follows from Proposition 11 of no. 6 and the definition of addition in \(\mathbf{P}(\mathbf{A})\) that \(\operatorname{cl}(\mathbf{M},\mathbf{N}) = \operatorname{cl}(\mathbf{M}) + \operatorname{cl}(\mathbf{N})\) for \(\mathbf{M}\) , \(\mathbf{N}\) in \(\mathfrak{J}\) , which shows that \(\operatorname{cl}\) is a homomorphism. If \(\mathbf{M} \in \mathfrak{J}\) is isomorphic to \(\mathbf{A}\) , there exists \(b \in \mathbf{B}\) such that \(\mathbf{M} = Ab\) and, as \(\mathbf{M}\) is invertible, there exists \(b' \in B\) such that \(b'b = 1\) , in other words \(b\) is invertible in \(\mathbf{B}\) ; the converse is immediate. Hence the kernel of \(\operatorname{cl}\) in \(\mathfrak{J}\) is \(\theta(\mathbf{B}^*)\) .

Let us now determine the image of cl. If \(M \in J\) , then \(M \otimes_{A} B = S^{-1}M = B\) (no. 5, Proposition 8 (c)), whence \(\operatorname{cl}(M) \in \operatorname{Ker}(\phi)\) . Conversely, let P be a projective A-module of rank 1 such that \(P_{(B)} = P \otimes_{A} B\) is B-isomorphic to B. As P is a flat A-module, the injection \(u: A \to B\) defines an injection \(u \otimes 1: P \to P_{(B)} = B\) and P is thus identified with a sub-A-module of B; by virtue of Proposition 8 (c) of no. 5, P is non-degenerate and Theorem 4 of no. 6 shows that P is invertible. The kernel of \(\phi\) is therefore equal to the image of \(\operatorname{cl}: J \to P(A)\) .

COROLLARY 1. For two invertible sub-A-modules of B to have the same image in C, it is necessary and sufficient that they be isomorphic.

COROLLARY 2. If the ring B is semi-local, the group C of classes of invertible sub-A-modules of B is canonically identified with the group P(A) of classes of projective A-modules of rank 1.

In this case \(\mathrm{P(B)} = 0\) (no. 3, Proposition 5).

Remark. The hypothesis of Corollary 2 is fulfilled in the two following cases:

\begin{enumerate}
\def\labelenumi{(\arabic{enumi})}
\tightlist
\item
  A is an integral domain and S is the set of elements \(\neq0\) of A, B then being the field of fractions of A. The invertible sub-A-modules of B are also called in this case invertible fractional ideals; those which are monogenous free A-modules Ab ( \(b\neq0\) in B) are just the fractional principal ideals defined in Algebra, Chapter VI, §1, no. 5.
\end{enumerate}

*(2) The ring A is Noetherian and S is the set of elements of A which are not divisors of 0 such that B is the total ring of fractions of A. In this case \(S = A - \bigcup_{i} p_i\) , where the \(p_i\) are the elements (finite in number) of Ass(A) (Chapter IV, §1), hence B is semi-local (§3, no. 5, Proposition 17). *


